% open-logic.tex
% open-logic-dev శైలితో మొత్తం పాఠ్యాన్ని నిర్మిస్తుంది

\documentclass[../include/open-logic-part]{subfiles}

\begin{document}

\clearpage

\begin{editorial}
ఈ ఫైలు ఓపెన్ లాజిక్ ప్రాజెక్ట్‌లో చేర్చిన మొత్తం విషయాన్ని లోడ్
చేస్తుంది. ఇలాంటివి చూపించినప్పుడు సంపాదకీయ గమనికలు, పాఠ్యాన్ని
ఎలా అందించాలో ఆలోచించకుండా ఫైలు నిర్మించబడిందని సూచిస్తాయి.
మీరు దీన్ని చదవగలుగుతున్నారంటే, ఈ PDFతో బోధించడం లేదా
దీనినే చదువుకోవడం బహుశా \emph{మంచిది కాదు}.

బోధనకు లేదా స్వయంగా అధ్యయనానికి మరింత అనుకూలమైన పాఠ్యాన్ని
తయారుచేసే అనేక విధానాలు ఓపెన్ లాజిక్ ప్రాజెక్ట్‌లో ఉన్నాయి.
ఉదాహరణకు, ముందస్తు అమరికలో అన్ని తార్కిక సంయోజకాలను
ప్రాథమికాలుగా తీసుకుని, నిరూపణల్లో ప్రతి సంయోజకానికి సంబంధించిన
అన్ని సందర్భాలనూ పూర్తి చేస్తుంది. కానీ కొన్ని సందర్భాలను
సాధనలుగా వదిలేయడం చాలా మెరుగైనది. ఓపెన్ లాజిక్ ప్రాజెక్ట్
పని ఇంకా కొనసాగుతోంది. సహకారాన్ని, మెరుగుదలను ప్రోత్సహించడానికి,
ముసాయిదా దశలోనే ఉన్న, సాధనలు ఇంకా లేని మొదలైన అంశాలనూ
ఇందులో చేర్చాం. ఒక పాఠ్యక్రమం కోసం నిర్మించిన PDFలో అలాంటి
విభాగాలను తీసివేస్తారు.

బోధన, అధ్యయనానికి మరింత అనుకూలమైన PDFల కోసం
\href{http://builds.openlogicproject.org/}{OLP వెబ్‌సైటులోని నమూనా పాఠ్యక్రమాలను}
చూడండి. మీ స్వంత పాఠ్యాన్ని తయారుచేయాలంటే,
\href{https://github.com/OpenLogicProject/OpenLogic/tree/master/courses/sample}{నమూనా నిర్వహణ ఫైలు}తో
మొదలుపెట్టవచ్చు; ఇంకా విస్తృతమైన ఉదాహరణల కోసం ఈ మూలం నుంచి
రూపొందించిన పాఠ్యపుస్తకాల మూలాలను చూడవచ్చు.
\end{editorial}

\olimport[sets-functions-relations]{sets-functions-relations-complete}

\olimport[propositional-logic]{propositional-logic}

\olimport[first-order-logic]{first-order-logic}

\olimport[model-theory]{model-theory}

\olimport[computability]{computability}

\olimport[turing-machines]{turing-machines}

\olimport[incompleteness]{incompleteness}

\olimport[second-order-logic]{second-order-logic}

\olimport[lambda-calculus]{lambda-calculus}

\olimport[many-valued-logic]{many-valued-logic}

\olimport[normal-modal-logic]{normal-modal-logic}

\olimport[applied-modal-logic]{applied-modal-logic}

\olimport[intuitionistic-logic]{intuitionistic-logic}

\olimport[counterfactuals]{counterfactuals}

\olimport[set-theory]{set-theory}

\olimport[methods]{methods}

\olimport[history]{history}

\olimport[reference]{reference}

\end{document}
